\documentclass[10pt,a4paper]{amsart}
\usepackage{amsmath,amssymb,amsthm,graphicx}
\usepackage{verbatim}
\newenvironment{reference}{\comment}{\endcomment}
%\newenvironment{reference}{}{}
\newenvironment{slogan}{\comment}{\endcomment}
\newenvironment{history}{\comment}{\endcomment}

% For commutative diagrams we use Xy-pic
\usepackage[all]{xy}

% We use 2cell for 2-commutative diagrams.
\xyoption{2cell}
% Theorem environments.
%
\theoremstyle{plain}
\newtheorem{theorem}[subsection]{Theorem}
\newtheorem{proposition}[subsection]{Proposition}
\newtheorem{lemma}[subsection]{Lemma}

\theoremstyle{definition}
\newtheorem{definition}[subsection]{Definition}
\newtheorem{example}[subsection]{Example}
\newtheorem{exercise}[subsection]{Exercise}
\newtheorem{situation}[subsection]{Situation}

\theoremstyle{remark}
\newtheorem{remark}[subsection]{Remark}
\newtheorem{remarks}[subsection]{Remarks}

\numberwithin{equation}{subsection}

% Macros
%
\def\lim{\mathop{\mathrm{lim}}\nolimits}
\def\colim{\mathop{\mathrm{colim}}\nolimits}
\def\Spec{\mathop{\mathrm{Spec}}}
\def\Hom{\mathop{\mathrm{Hom}}\nolimits}
\def\Ext{\mathop{\mathrm{Ext}}\nolimits}
\def\SheafHom{\mathop{\mathcal{H}\!\mathit{om}}\nolimits}
\def\SheafExt{\mathop{\mathcal{E}\!\mathit{xt}}\nolimits}
\def\Sch{\mathit{Sch}}
\def\Mor{\mathop{\mathrm{Mor}}\nolimits}
\def\Ob{\mathop{\mathrm{Ob}}\nolimits}
\def\Sh{\mathop{\mathit{Sh}}\nolimits}
\def\NL{\mathop{N\!L}\nolimits}
\def\CH{\mathop{\mathrm{CH}}\nolimits}
\def\proetale{{pro\text{-}\acute{e}tale}}
\def\etale{{\acute{e}tale}}
\def\QCoh{\mathit{QCoh}}
\def\Ker{\mathop{\mathrm{Ker}}}
\def\Im{\mathop{\mathrm{Im}}}
\def\Coker{\mathop{\mathrm{Coker}}}
\def\Coim{\mathop{\mathrm{Coim}}}

% Boxtimes
%
\DeclareMathSymbol{\boxtimes}{\mathbin}{AMSa}{"02}

%
% Macros for moduli stacks/spaces
%
\def\QCohstack{\mathcal{QC}\!\mathit{oh}}
\def\Cohstack{\mathcal{C}\!\mathit{oh}}
\def\Spacesstack{\mathcal{S}\!\mathit{paces}}
\def\Quotfunctor{\mathrm{Quot}}
\def\Hilbfunctor{\mathrm{Hilb}}
\def\Curvesstack{\mathcal{C}\!\mathit{urves}}
\def\Polarizedstack{\mathcal{P}\!\mathit{olarized}}
\def\Complexesstack{\mathcal{C}\!\mathit{omplexes}}
% \Pic is the operator that assigns to X its picard group, usage \Pic(X)
% \Picardstack_{X/B} denotes the Picard stack of X over B
% \Picardfunctor_{X/B} denotes the Picard functor of X over B
\def\Pic{\mathop{\mathrm{Pic}}\nolimits}
\def\Picardstack{\mathcal{P}\!\mathit{ic}}
\def\Picardfunctor{\mathrm{Pic}}
\def\Deformationcategory{\mathcal{D}\!\mathit{ef}}
\usepackage{cleveref}
\crefname{originaldefinition}{original definition}{original definitions}
\Crefname{originaldefinition}{Original definition}{Original definitions}
\makeatletter
\newlabel{definition-quasi-compact}{{03HE}{source}{}{sourcepointer-1}{}}
\newlabel{definition-quasi-compact@cref}{{[originaldefinition][0][]03HE}{[1][0][]source}{}{}{}}
\makeatother
\crefname{originallemma}{original lemma}{original lemmas}
\Crefname{originallemma}{Original lemma}{Original lemmas}
\makeatletter
\newlabel{spaces-lemma-base-change-immersions}{{02YW}{source}{}{sourcepointer-2}{}}
\newlabel{spaces-lemma-base-change-immersions@cref}{{[originallemma][0][]02YW}{[1][0][]source}{}{}{}}
\makeatother
\crefname{originallemma}{original lemma}{original lemmas}
\Crefname{originallemma}{Original lemma}{Original lemmas}
\makeatletter
\newlabel{spaces-properties-lemma-open-subspaces}{{03BZ}{source}{}{sourcepointer-3}{}}
\newlabel{spaces-properties-lemma-open-subspaces@cref}{{[originallemma][0][]03BZ}{[1][0][]source}{}{}{}}
\makeatother
\crefname{originallemma}{original lemma}{original lemmas}
\Crefname{originallemma}{Original lemma}{Original lemmas}
\makeatletter
\newlabel{spaces-properties-lemma-points-cartesian}{{03H4}{source}{}{sourcepointer-4}{}}
\newlabel{spaces-properties-lemma-points-cartesian@cref}{{[originallemma][0][]03H4}{[1][0][]source}{}{}{}}
\makeatother
\crefname{originallemma}{original lemma}{original lemmas}
\Crefname{originallemma}{Original lemma}{Original lemmas}
\makeatletter
\newlabel{spaces-properties-lemma-quasi-compact-space}{{03E4}{source}{}{sourcepointer-5}{}}
\newlabel{spaces-properties-lemma-quasi-compact-space@cref}{{[originallemma][0][]03E4}{[1][0][]source}{}{}{}}
\makeatother
\title{A quasi-compact morphism of algebraic spaces induces a quasi-compact map of topological spaces}
\author{The Stacks Project Authors}
\date{stacks-project:spaces-morphisms.tex:local-snapshot-20261012; GNU FDL 1.2 or later}
\begin{document}
\maketitle
\setcounter{subsection}{1}
\noindent\textit{Editorial note.} All delivered mathematical text is the author's own, copied from the pinned Stacks Project chapter file spaces-morphisms.tex; no statement or proof step was rewritten, reordered or removed.\par\medskip
\noindent\textit{Editorial note.} The proof is a topological translation of the definition of quasi-compactness of a morphism; the description of the points of the fibre product is what identifies the preimage.\par\medskip
\noindent\textit{Editorial note.} The five references of the counted spans are recorded as non-load-bearing navigation pointers; the argument that consumes them stays inside the counted proof span.\par\medskip
\noindent\textit{Editorial note.} Preamble: the three native excerpt spans of sources/stacks-spaces-morphisms/preamble.tex (lines 8-13, 14-18 and 154-217). The chapter file itself carries no preamble text.\par\medskip
\noindent\textit{Allowed standard prerequisites (editorial).}\begin{itemize}
\item Every open subset of |Y| is the underlying set of an open subspace. Properties of Spaces, Lemma spaces-properties-lemma-open-subspaces, used at proof lines 1274-1276. Turns the quasi-compact open V into an open subspace Y'.
\item Quasi-compactness of an algebraic space tested on the underlying topological space. Properties of Spaces, Lemma spaces-properties-lemma-quasi-compact-space, used at proof lines 1277-1279 and 1285-1287. Converts topological quasi-compactness into quasi-compactness of Y' and back for |X'|.
\item The definition of quasi-compactness for a morphism of algebraic spaces. Definition definition-quasi-compact, used at proof lines 1279-1280. Makes X' = Y' x\_Y X quasi-compact.
\item Base change of open immersions and the description of points of a fibre product. Spaces, Lemma spaces-lemma-base-change-immersions and Properties of Spaces, Lemma spaces-properties-lemma-points-cartesian, used at proof lines 1281-1284. Identifies X' as an open subspace of X and |X'| with the preimage of V.
\end{itemize}\medskip
\ifdefined\hypertarget\hypertarget{sourcepointer-1}{}\fi
\noindent\textit{Source reference.} Original definition 03HE: focused-seeds/AGP/sources/stacks-snapshot-20261011/tags.txt, source lines 4563--4563. Navigation pointer into the Stacks Project; the standard fact it supplies is applied inside the counted proof span and the argument resting on it stays there.\par\medskip
\ifdefined\hypertarget\hypertarget{sourcepointer-2}{}\fi
\noindent\textit{Source reference.} Original lemma 02YW: focused-seeds/AGP/sources/stacks-snapshot-20261011/tags.txt, source lines 3891--3891. Navigation pointer into the Stacks Project; the standard fact it supplies is applied inside the counted proof span and the argument resting on it stays there.\par\medskip
\ifdefined\hypertarget\hypertarget{sourcepointer-3}{}\fi
\noindent\textit{Source reference.} Original lemma 03BZ: focused-seeds/AGP/sources/stacks-snapshot-20261011/tags.txt, source lines 4366--4366. Navigation pointer into the Stacks Project; the standard fact it supplies is applied inside the counted proof span and the argument resting on it stays there.\par\medskip
\ifdefined\hypertarget\hypertarget{sourcepointer-4}{}\fi
\noindent\textit{Source reference.} Original lemma 03H4: focused-seeds/AGP/sources/stacks-snapshot-20261011/tags.txt, source lines 4553--4553. Navigation pointer into the Stacks Project; the standard fact it supplies is applied inside the counted proof span and the argument resting on it stays there.\par\medskip
\ifdefined\hypertarget\hypertarget{sourcepointer-5}{}\fi
\noindent\textit{Source reference.} Original lemma 03E4: focused-seeds/AGP/sources/stacks-snapshot-20261011/tags.txt, source lines 4443--4443. Navigation pointer into the Stacks Project; the standard fact it supplies is applied inside the counted proof span and the argument resting on it stays there.\par\medskip
\begin{lemma}
\label{lemma-quasi-compact-is-quasi-compact}
Let $S$ be a scheme. If $f : X \to Y$ is a quasi-compact morphism of
algebraic spaces over $S$, then the underlying map
$|f| : |X| \to |Y|$ of topological space is quasi-compact.
\end{lemma}

\begin{proof}
Let $V \subset |Y|$ be quasi-compact open. By Properties of Spaces,
Lemma \ref{spaces-properties-lemma-open-subspaces}
there is an open subspace $Y' \subset Y$ with $V = |Y'|$.
Then $Y'$ is a quasi-compact algebraic space by
Properties of Spaces, Lemma \ref{spaces-properties-lemma-quasi-compact-space}
and hence $X' = Y' \times_Y X$ is a quasi-compact algebraic
space by Definition \ref{definition-quasi-compact}.
On the other hand, $X' \subset X$ is an open subspace
(Spaces, Lemma \ref{spaces-lemma-base-change-immersions})
and $|X'| = |f|^{-1}(|X'|) = |f|^{-1}(V)$ by
Properties of Spaces, Lemma \ref{spaces-properties-lemma-points-cartesian}.
We conclude using
Properties of Spaces, Lemma \ref{spaces-properties-lemma-quasi-compact-space}
again that $|X'|$ is a quasi-compact open of $|X|$ as desired.
\end{proof}

\end{document}
